% Numbers for the paper. Pre-freeze values are filled in by hand from results/ and prereg/. The results texts were
% written by hand from results/ (main), results/heldout and results/audit after the runs finished on 2026-11-06;
% make_results.py only fills macros that still print TBD. Tables come from results_tables.tex (generated).
\newcommand{\resAbstractFinding}{that C verifies far more repairs than B for both primary models (Qwen3.5-2B: 27.7\%
of families vs 2.2\%; Qwen3.5-4B: 33.0\% vs 17.6\%; Holm $p<0.001$), and at least as many as BP (2B: +5.0
percentage points, $p_\text{Holm}=0.012$; 4B: a tie), while spending about 43\% fewer model tokens per verified
repair than BP and never asking the model for a proof. The held-out split replicates C over B and shows a C/BP tie
for both models. The gain comes from the automation more than from forbidding proofs: proof-writing agents that are
given the same portfolio first close most of the gap, and almost all of their successes are the portfolio's proofs,
not their own}

\newcommand{\resRQOneText}{\Cref{tab:rq1} gives the four confirmatory tests on the 106 \texttt{main} families that
carry a bug in the primary pool, and \cref{fig:rates} shows the rates for every condition.
\textbf{C vs B} is large and significant for both primary models: Qwen3.5-2B 27.7\% vs 2.2\% (+25.5\,pp, 95\% CI
[+18.2, +33.0]); Qwen3.5-4B 33.0\% vs 17.6\% (+15.4\,pp, [+9.4, +21.7]); $p_\text{Holm}<0.001$ for both. C
repairs 28 families that B never repairs for 2B (B: 0 that C misses), and 19 against 1 for 4B.
\textbf{C vs BP} is small. For 2B, C is ahead by +5.0\,pp ([+1.6, +8.8], $p_\text{Holm}=0.012$; 10 families only
by C, 2 only by BP). For 4B the two are tied: 33.0\% vs 32.7\% (+0.3\,pp, [$-$2.8, +3.5], $p_\text{Holm}=0.50$; 3
vs 2 families). The pre-freeze power analysis predicted exactly this: C vs BP had 0.71 power for 2B and 0.08 for
4B at the available families (\cref{tab:power}), so the 4B tie is reported as a tie, not as evidence of no
difference. Per-run rates (\cref{tab:rates}) and outcome labels (\cref{tab:outcomes}) tell the same story.
No run in any condition was \emph{verified but behaviourally wrong}, the outcome that would expose a weak
specification. Most failures in the proof-writing conditions are submissions rejected by the static checks;
they are mostly malformed rather than adversarial (lemmas placed in the helper region, a second definition of
the function, or proof text in the code region), and they are far rarer in C and C-plain (1,690 and 1,586
rejected attempts) than in A, B and BP (7,187, 5,662 and 5,764).}

\newcommand{\resCvsBPText}{The 4B tie stands as the result; four further facts say what it means. (i)~\emph{Power.} With 106 families, the pre-freeze simulation gave C vs BP a power of 0.08 for
4B (\cref{tab:power}): a difference of a few points could not have been detected, so the tie says ``no large
difference'', not ``equal''. (ii)~\emph{Where BP's successes come from.} BP is B with the portfolio run first. 180
of its 182 verified runs for 4B, and all 110 for 2B, were closed by the portfolio, not by the model's proof; the
proof-writing half of BP added two repairs out of 693 runs. BP is, in effect, C plus an expensive and almost
always failing proof attempt. (iii)~\emph{Cost.} For the same success rate, C spends 37k model tokens per verified
repair against BP's 65k (43\% fewer) and 68\,s of inference per run against 162\,s, because the model writes no
proof text; for 2B, 50k against 89k. (iv)~\emph{Integrity.} BP's submissions are rejected by the static checks
3.4 times as often as C's (5,764 against 1,690 rejected attempts), because proof text is where most malformed
submissions arise. The practical reading is that once automation is available, asking the model for a proof adds
cost and failure modes but no verified repairs.}

\newcommand{\resRQTwoText}{We classify each success as a \emph{revert} (identical to the reference, or within the
pre-registered distance of a reference that the portfolio proves), \emph{shaping} (the reference itself is
\emph{not} provable by the portfolio, so the verified code must have a different, automation-friendly shape) or
other. Shaping is common: 42 of 150 C successes for 2B and 81 of 191 for 4B are shaping, against 72 and 56
reverts. Where the reference is provable by the portfolio, C and BP both succeed on most runs (4B: 89.4\% and
91.1\%; B: 36.6\%). Where it is not, the rates fall for every condition, but C stays ahead of B (4B: 14.2\% vs
6.0\%; 2B: 7.4\% vs 0.0\%) and level with BP (12.3\%). Restricting the paired family-level test to the 85
\texttt{main} families whose reference the portfolio cannot prove, so that no success can be a plain revert
(\cref{tab:robust}, exploratory), C still beats B for every model (2B +10.6\,pp, 4B +7.8\,pp, Gemma +15.3\,pp;
$p\le0.002$), and is ahead of BP by +4.3 (2B, $p=0.035$), +1.2 (4B, $p=0.34$) and +7.1\,pp (Gemma, $p=0.055$). The
headline is therefore not carried by easy cases alone, although rates on these families are about half the
overall ones. Fifteen real before/after cases, each with the bug, the reference, C's verified code, the closing
tactic and what B did on the same run, are in \texttt{docs/code\_shapes.md}; \cref{fig:shape} shows one.
\textbf{Study S} asks the harder question directly:
starting from a \emph{correct} reference that the portfolio cannot prove, can a model re-shape it so that the
portfolio can? Rarely, and the structured hint does not help: 3.7\% (S) vs 4.1\% (S-plain) of runs for 2B and
12.7\% vs 12.7\% for 4B (family-level differences $-$0.4 and 0.0\,pp, $p=0.81$ and $0.60$). This suggests that
shaping in the repair setting happens mostly when fixing the bug and re-shaping the code coincide, not as a skill
of its own.}

\newcommand{\resRQThreeText}{The structured hint (witness, best unsolved goal, structural hint) does not beat a plain
retry. C vs C-plain: 2B 27.7\% vs 28.9\% ($-$1.3\,pp, [$-$3.8, +1.3], $p=0.89$), 4B 33.0\% vs 32.7\% (+0.3\,pp,
[$-$2.2, +2.8], $p=0.50$); 3 vs 3 and 2 vs 1 discordant families. C-plain is slightly cheaper (4B: 9.1k vs 10.3k
tokens and 58 vs 68\,s of inference per run). The benefit of C comes from the code-only contract and the
portfolio, not from the feedback text.}

\newcommand{\resRQFourText}{Recursion is the dominant blocker. Among the 184 eligible references (condition T), the
portfolio closes none of the 101 that use any recursion and 49.4\% of the rest; \texttt{let rec} (0 of 77) and
\texttt{match} (1.5\% vs 33.9\%) follow. The same holds for code the models write: among test-passing candidates
from the code-only conditions, 0.7\% of those with recursion are closed against 67.2\% without, 10.3\% with
\texttt{match} against 62.6\%, 29.7\% with \texttt{a[i]!} indexing against 58.7\%, and 23.7\% with
\texttt{Nat}/\texttt{Int} casts against 58.1\%. Folds and list combinators cost little ($-$4 and $-$10\,pp).
Verified repairs are closed mostly by the cheapest entries: for C, \texttt{simp} closes 146 of 399, \texttt{unfold}
+ \texttt{omega} 90, \texttt{simp\_all} 58, \texttt{simp}+\texttt{omega} 38 and \texttt{simp}+\texttt{grind} 36; the
induction entries never close a repair. Automation-friendly code, in practice, means non-recursive code built from
library combinators.}

\newcommand{\resRQFiveText}{\textbf{Cost.} C is the cheapest route to a verified repair (\cref{fig:costs}). For
Qwen3.5-4B, C spends 37k model tokens per verified repair, against 65k for BP and 165k for B, and 68\,s of
inference per run against 162\,s and 150\,s; for 2B, 50k against 89k (BP) and 1.5M (B). C spends more Lean time
per run (18\,s vs 6\,s for B with 4B) because the portfolio runs on every candidate. Truncating the generation
budget to 2 or 4 shows that C's lead is already present at $N=2$ (4B: 21.8\% vs 6.8\% of runs for B). Patch sizes
(normalised edit distance from the buggy start, successes only) are similar across conditions (4B: C 0.67, BP
0.52, B 0.77). \textbf{Scale.} Going from 2B to 4B helps B most (2.2\% $\to$ 17.6\%) and C least (27.7\% $\to$
33.0\%): the smaller model gains most from not writing proofs. \textbf{Second family.} Gemma-4-E4B (secondary)
repeats the pattern: C 30.2\% vs B 1.9\% (+28.3\,pp, $p<0.001$) and C vs BP 30.2\% vs 27.4\% (+2.8\,pp, [$-$3.8,
+9.4], $p=0.29$). \textbf{Proof specialist.} Goedel-Code-Prover-8B, run as a reference in A and B only, verifies
16.5\% and 20.8\% of runs, below Qwen3.5-4B in C (27.6\%), which has half the parameters and no proof training.
\textbf{Energy} (estimate, deviation D-4; \cref{tab:energy}). The study's energy counter read the wrong GPU, so
we calibrated joules per prompt and completion token for each model on our slice after the study. Each model was
replayed on prompts it had received, at the study's serving settings. C is also the most energy-efficient
condition: for Qwen3.5-4B it uses 2.6\,kJ per verified repair [2.4, 2.9], against 5.9\,kJ for BP and 13.4\,kJ for
B; for 2B, 2.0\,kJ against 3.6\,kJ and 56\,kJ; for Gemma, 2.6\,kJ against 8.8\,kJ and 232\,kJ. These are 0.4--0.5\,g of CO$_2$
per repair in C. The held-out split gives the same ordering (4B: 3.2, 7.4 and 24\,kJ). C's extra Lean time costs
less than the model tokens it saves. All study runs together used an estimated 3.1\,kWh (2.3\,kg CO$_2$). Running two
jobs side by side may raise absolute figures by up to about 1.6 times, but equally for every condition.}

\newcommand{\resSecondaryText}{\textbf{B without automation.} Many of B's own proofs are a single automation call.
Counting those as failures (B\_noauto) lowers B from 17.6\% to 12.3\% for 4B (C vs B\_noauto +20.8\,pp); for 2B,
B falls to 0.9\%. \textbf{Bhammer}, which tells B the portfolio's tactics in the prompt, helps B a little (4B:
24.2\% vs 17.6\%, +6.6\,pp, [+1.3, +11.9], $p=0.011$; 2B: 5.3\% vs 2.2\%, $p=0.057$) but stays below C (4B:
$-$8.8\,pp, $p=0.001$; 2B: $-$22.3\,pp, $p<0.001$); 95 of its 131 successes for 4B are automation-only proofs. Knowing
the tactics is not the same as having them run. \textbf{B$\to$C cross-evaluation.} We re-verified every B attempt
whose code compiled but whose proof failed through C's code-only pipeline. C's pipeline would have verified
30.8\% of these attempts for 4B (797 of 2,588), 34.0\% for 2B and 40.7\% for Gemma, and would have rescued 17.3\%,
16.8\% and 23.1\% of B's failed runs. None of these candidates passed the proof but failed the tests, so the
cross-evaluation found no weak specification. In B, the code was often right and the proof was the bottleneck.
\textbf{Natural bugs only.} Natural bugs cover only 33 \texttt{main} families, which is why the pre-registered
fallback pool applies. On those families alone (\cref{tab:robust}, exploratory), C beats B for every model (2B
14.1\% vs 0.0\%, $p=0.031$; 4B 18.2\% vs 0.0\%, $p=0.008$; Gemma 18.2\% vs 0.0\%, $p=0.016$), and C is ahead of BP
by 5--9\,pp without significance ($p$ from 0.06 to 0.13). Mutants alone reproduce the pooled results. Natural bugs
are harder for every condition, and the small pool limits what can be said about them beyond the direction.
\textbf{Re-verification.} \resReverifyText
\textbf{Audit.} The 10\% stratified audit plus every flagged success (333 runs, 126 distinct programs) found no
specification gaming in an automated pre-screen; the human audit is pending (deviation D-3). 27 runs decide the
postcondition proposition directly or carry dead code, and are reported as unnatural but correct. The spec-copying
flag fired on 239 of the audited runs, all of which restate an exact, executable postcondition.}

\newcommand{\resHeldoutText}{The held-out split (46 families with a bug in the primary pool) was run once and
analysed separately (\cref{tab:heldout}). C vs B replicates for both models: 2B 20.3\% vs 1.4\% (+18.8\,pp,
[+9.4, +29.0]), 4B 29.7\% vs 10.1\% (+19.6\,pp, [+10.1, +30.4]); $p_\text{Holm}=0.001$ for both. C vs BP is a tie for
both models (20.3\% vs 20.3\% and 29.7\% vs 29.7\%; $p_\text{Holm}=1.00$), so the 2B advantage of C over BP on
\texttt{main} does not replicate. No condition repaired any of the 10 held-out families with natural bugs.}

\newcommand{\resConclusionText}{Across 12,249 study runs on one GPU slice, letting automation carry the proof is the
most effective and cheapest way for small models to produce end-to-end verified repairs. C beats proof writing
(B) decisively, on \texttt{main} and on the held-out split, for two Qwen sizes and for Gemma. Against a
proof-writing agent that gets the same automation first (BP), C ties or wins narrowly, and the narrow win does not
replicate on held-out. It is clearly cheaper: about 43\% fewer tokens per verified repair and less than half the
inference time with 4B (and, by a post-hoc estimate, less than half the energy). The structured hint adds nothing
over a plain retry; recursion, \texttt{match} and array
indexing are what defeat the automation; and models rarely re-shape correct but unprovable code. The practical
lesson is to give small models a fixed automation portfolio and ask them only for code. Further gains will come
from widening what automation can prove, especially recursion, rather than from better feedback text.}

% from results/reverify/summary.md (scripts/reverify.py, 2026-11-14)
\newcommand{\resReverifyText}{After the study, every one of the 1,858 logged successes (all study folders) was given
again to the trusted judge from its logged code and proof, with the cache disabled and on a different machine and
operating system (a Windows laptop, same Lean and Mathlib). All 1,858 re-verify, each with the same proof source and
the same closing portfolio entry as logged (D-7). The verdicts therefore do not rest on the logs alone.}

\newcommand{\resFreezeDigest}{\texttt{3f2744e1\ldots}, 19 files, 2026-11-05}
\newcommand{\resTCoverage}{22.3\%: basic 35.6\%, advanced 5.0\%}
\newcommand{\resNaturalFamilies}{33}
\newcommand{\resPowerTable}{%
\begin{tabular}{@{}llrrr@{}}
\toprule
Comparison & Model & Point & Pessimistic & Families needed \\
\midrule
C vs B  & Qwen3.5-2B & 1.000 & 0.662 & 34 \\
C vs B  & Qwen3.5-4B & 0.996 & 0.128 & 57 \\
C vs BP & Qwen3.5-2B & 0.706 & 0.000 & 136 \\
C vs BP & Qwen3.5-4B & 0.082 & 0.000 & not reachable \\
\bottomrule
\end{tabular}}
